\subsection{Security of LLM-Assisted Code Generation}

\noindent \textbf{Poison Detection in LLM Code Generation.}
Recent studies investigate malicious behavior during inference 
for code-generation assistants.
Prior works\cite{zeng2025inducing,yi2025benchmarking} show that vulnerable code generation can 
be induced through malicious external content or tool-mediated 
contexts.
These findings motivate systematic assessment of prompt-conditioned 
attack surfaces in practical developer workflows.
In contrast, this thesis demonstrates a more fundamental 
problem that does not require the LLM to 
access any external sources during inference. 


\noindent \textbf{Poisoning in LLM Training Pipelines.}
Data poisoning has emerged as a major threat to modern ML systems, 
with early analyses in computer vision and broader surveys of poisoning 
defenses\cite{cina2023wild,raghavan2022improved,goldblum2022dataset}.
For LLMs specifically, web-scale poisoning feasibility has been 
demonstrated\cite{carlini2024poisoning},
with subsequent studies analyzing poisoning across pretraining and 
fine-tuning regimes\cite{zhao2025data,jiang2024turning}.
This thesis differs by (i) targeting malicious code 
generation, which poses immediate execution risks, rather 
than natural-language outputs, and (ii) auditing production 
LLMs for evidence of existing, passive poisoning in 
their training corpora, rather than new inference-time active attacks.



\noindent \textbf{Security Benchmarks for LLM-Generated Code.}
Another line of work builds benchmark suites for security 
evaluation of generated code,
including SecCodePLT, SafeGenBench, and 
CodeLMSec\cite{yang2024seccodeplt,li2025safegenbench,debenedetti2023a}.
These benchmarks mainly emphasize CWE-oriented vulnerabilities,
while this thesis additionally focuses on harmful 
artifact reproduction in generated code (e.g., malicious URLs and endpoints).
